\chapter{Interface de rede}

\chapterauthor{Felipe Pregnolatto Melo, Raul Yoshiyuki Koma, Vicente Pinto Tomás Junior}

%-------------------------------------------------------------------------------
\section{Descrição}
%-------------------------------------------------------------------------------
% Descreva em alto nível o que este módulo faz e qual problema ele resolve
% dentro do sistema operacional.
%

A placa de rede é a responsável pelo envio e recebimento de pacotes para a rede em baixo nível. Por ser um componente externo, é responsabilidade do sistema operacional fazer o acordo e comunicação em memória com o \emph{driver}.

A interface de rede então é o sistema que faz essa comunicação, estabelece os acordos, define as posições de memória onde os pacotes chegam e faz o tratamento da leitura e envio dos pacotes para a internet.



%-------------------------------------------------------------------------------
\section{Decisões de Arquitetura e Design}
%-------------------------------------------------------------------------------
% Justificar o porque de ter feito as coisas de uma certa maneira.
%
% Tópicos a cobrir:
% - Quais alternativas foram consideradas?
% - Por que a abordagem escolhida foi a melhor para este projeto?
% - Quais são as limitações ou trade-offs da abordagem escolhida?
%

Como a placa de rede é um driver externo, em um kernel real ele é uma peça de \emph{hardware} adicional à maquina. Por sorte o QEMU ARM contém uma série de disposítivos simulados que podem ser adicionados e configurados, chamados série Virtio (\emph{Virtual I/O Device}), sendo um deles um disposítivo de rede.

O disposítivo escolhido foi o disposítivo de rede versão 1.0 da série de disposítivos Virtio de rede, uma versão já legada, o motivo dessa escolha foi simplesmente por ser a única versão que era conectada pelo QEMU, o mótivo disso ocorrer não foi descoberto pelo grupo.

As limitações então são responsabilidade do sistema esperar o disposítivo versão 1.0, não podendo ser iniciado caso a versão seja diferente. \cite{virtio:networkdevice}



%-------------------------------------------------------------------------------
\section{Detalhes da Implementação}
%-------------------------------------------------------------------------------
% Descreva como a solução foi implementada, conectando o código
% à teoria.
%
% Tópicos a cobrir:
% - Principais estruturas de dados (structs) e suas funções.
% - Algoritmos chave utilizados.
% - Fluxo de execução das funções mais importantes.
% - Inclua pequenos trechos de código com o pacote 'listings'.
%

\subsection{Arquivos}

A implementação foi dividida nos seguintes arquivos:

\begin{description}
    \item[\texttt{virtio.h}] Cabeçalho que define a struct de um socket
    \begin{lstlisting}[language=C]
        struct socket {
            uint32_t em_uso;
            uint16_t porta_local;
            uint8_t rx_buffer[4096];
            uint32_t rx_tamanho;
            struct process *processo_bloqueado;
        };
    \end{lstlisting}
    Além de definir as funções da API.

    \item[\texttt{protocol\_headers.h}] Cabeçalho que define várias structs dos cabeçalhos de diversos protocolos para um pacote da internet, como IPV4, ARP e UDP.

    \item[\texttt{main.c}] Onde é executado a função de inicialização da placa de rede.

    \item[\texttt{virtio.c}] Arquivo principal de implementação e as funções internas e externas .
\end{description}

\subsection{Fluxo de execução}

    Antes de explicar o fluxo, é necessário explicar a estrutura do que são virtqueues, o modelo de dados usado pela placa de rede.

\subsubsection{VirtQueues}

\begin{lstlisting}[language=C]
struct virtq_desc {
    uint64_t addr;
    uint32_t len;
    uint16_t flags;
    uint16_t next;
} __attribute__((packed));

struct virtq_avail {
    uint16_t flags;
    uint16_t idx;
    uint16_t ring[];
} __attribute__((packed));

struct virtq_used_elem {
    uint32_t id;
    uint32_t len;
} __attribute__((packed));

struct virtq_used {
    uint16_t flags;
    uint16_t idx;
    struct virtq_used_elem ring[];
} __attribute__((packed));


    /* Estrutura para o kernel controlar o estado da fila */
struct virtqueue_info {
    struct virtq_desc *desc;
    struct virtq_avail *avail;
    struct virtq_used *used;
    uint16_t last_used_idx; // Controla até onde o kernel já leu
    uint16_t avail_idx;     // Controla onde colocar o próximo buffer 
};

// Filas globais (0 = RX, 1 = TX)
struct virtqueue_info virtqueues[2];
\end{lstlisting}

Virtqueues tem 3 partes importantes para entender sua estrutura

A parte descritora, que descreve os dados nos buffers, como seu endereço, seu tamanho e seu tipo (se é envio ou recebimento).

O anel de pacotes avaliaveis, onde o kernel escreve os pacotes que devem ser enviados pela placa de rede quando tiver oportunidade, tendo estrutura de uma fila circular.

O anel usado, onde o kernel escreve os pacotes recebidos pela placa, tendo estrutura de uma fila circular.


\subsubsection{Inicialização}

A função \texttt{virtio\_net\_init()} é responsável pela inicialização e ela segue o seguinte fluxo:

\begin{enumerate}
\item Marca o endereço da placa como \texttt{0x0A003e00}
\item Verifica se o endereço corresponde a um dispositivo virtio e se o dispositivo é um de rede da versão 1.0.
\item Começa a registrar no registrador de status, setando como 0.
\item Marca o bit 1 para dizer que o kernel reconhece o dispositivo
\item Marca o bit 2 para dizer que o kernel sabe como utilizar o dispositivo e dá tamanho da página em bytes.
\item Configura as virtqueues, inicializando a primeira fila como recebimento e a segunda como de envio.
\item Marca o bit 3 para dizer que o kernel terminou de configurar o dispositivo e ele esta pronto para uso.
\item Habilita a função de interrupção.
\item Inicializa os sockets com hashtable.
\item Envia um ARP para avisar o QEMU o endereço da placa.
\end{enumerate}

\subsubsection{Tratador de interrupção}

Uma interrupção na rede significa que um pacote chegou de fora, então a interrupção deve ser enviada para que em algum momento o pacote seja lido e entregue aos processos. Mesmo com a versão atual não tratando interrupções, a função ainda é usada como padrão de leitura.

\begin{lstlisting}[language=C]
while(rx->last_used_idx != rx->used->idx){

    ... faz a leitura dos pacotes baseado em portas

uint16_t novo_avail_idx = rx->avail_idx % VIRTQ_SIZE;
rx->avail->ring[novo_avail_idx] = desc_id;
rx->avail_idx++;
                
rx->last_used_idx++;
}
\end{lstlisting}

A função de interrupção então quando chamada, percorre toda a lista circular do anel de espaço usados da virqueue e se o pacote é um udp, envia os dados para o processo que abriu o socket do destino.

\subsubsection{Envio de pacotes}


sys\_recv(int32\_t port, void *buffer\_user, uint32\_t tam\_max)
Utiliza o método de pooling para chamar o tratamento de interrupção, então lê os dados do socket e retorna a quantidade de dados retornados enquanto escreve no buffer do usuario.

A função de envio de pacotes não está corretamente implementada.





%-------------------------------------------------------------------------------
\section{Interface Pública (API)}
%-------------------------------------------------------------------------------
% Documente como outros módulos do kernel podem usar este módulo.
% Essencialmente, é a documentação do arquivo .h correspondente.
%
% Para cada função pública:
% - Nome da função.
% - Descrição do que ela faz.
% - Parâmetros (o que cada um significa).
% - Valor de retorno (o que ela retorna e em que casos).
%

\subsubsection{virtio\_net\_init}

\textbf{Descrição:}
Inicializa a o dispositivo de rede.

\textbf{Parâmetros:}
Nenhum

\textbf{Retorno:}
Retorna 1 caso a inicialização tenha sucessido, -1 caso tenha falho e 2 caso a placa tenha sido configurada mas tenha ocorrido alguma falha na configuração da função de interrupção.


\subsubsection{socket}

\textbf{Descrição:}
Abre um socket com a porta indicada. As portas disponíveis no sistema são as 5000 até 5015 e a 5555.

\textbf{Parâmetros:}
\begin{itemize}
\item \texttt{porta}: Número da porta.
\end{itemize}

\textbf{Retorno:}
Retorna o ID do socket para uso nas outras funções, retorna -1 caso não tenha aberto nenhum.

\subsubsection{recv}

\textbf{Descrição:}
Lê os pacotes recebidos pelo kernel.

\textbf{Parâmetros:}
\begin{itemize}
\item \texttt{port}: Número da porta de recebimento.
\item \texttt{buffer\_user}: Endereço do buffer onde a função vai escrever os conteudos dos pacotes enviados.
\item \texttt{tam\_max}: Tamanho do máximo de bytes que serão escritos no buffer.
\end{itemize}

\textbf{Retorno:}
Retorna a quantidade de bytes lidos.


\subsubsection{send}
Não implementado.



%-------------------------------------------------------------------------------
\section{Testes e Verificação}
%-------------------------------------------------------------------------------
% Mostre como o funcionamento do módulo foi testado e validado.
%
% Tópicos a cobrir:
% - Descreva a rotina de teste que foi adicionada à função kmain().
% - Inclua a saída esperada (e obtida) na console serial do QEMU.
% - Explique quais casos de borda foram testados.
%
No Kmain(), é executado apenas a função de inicialização do dispositivo de rede. É colocado em uma variável o valor de retorno e printa a mensagem "Placa de rede iniciada." caso o retorno seja 1. O resultado obtido quando a inicialização é feita pelo make run-net é essa.

O resto dos testes não são feitos na aba de programas do shell, onde existem os programas \texttt{rede1} e \texttt{rede2}. \texttt{rede2} não está funcionando por utilizar do send.
O programa \texttt{rede1} pede uma porta de argumento e retorna o payload do pacote recebido nessa porta se existe algo. Essa funcionalidade pode ser testada dando um netcat udp para o endereço local e porta requisitada antes de dar o comando rede1. Como por exemplo:


% Nao sei como colocar o simbolo do |
- echo "Teste de Rede UFSKernel" {pipe} nc -u -w1 127.0.0.1 5555

Ao fazer esse echo antes de dar o comando pela primeira vez, o retorno é sempre correto e ele retorna o texto do echo completo.

Os pacotes podem ser visualizados após a execução do kernel também com o \texttt{qemu\-trafego.pcap} que sempre é gerado. Ele pode ser lido com o programa tshark pelo terminal com o comando:

- tshark -r qemu-trafego.pcap 

Essa leitura sempre retorna os pacotes recebidos e o ARP inicial corretamente.



